% pp2lp contribution for ICSPA / AFADL 2026
% Standalone buildable: lualatex afadl.tex
\documentclass[11pt]{article}
\usepackage{amsmath,amssymb,mathtools}
\usepackage{ifluatex}
\ifluatex\usepackage{fontspec}\fi
\usepackage{xcolor}
\usepackage{tikz}
\usetikzlibrary{positioning, calc}
\usepackage{listings}
\usepackage{tcolorbox}
\tcbuselibrary{listings, skins}
\usepackage{booktabs}
\usepackage[margin=2.5cm]{geometry}
\usepackage{hyperref}

% ── Colours ───────────────────────────────────────────────────────
\definecolor{lpKeyword}{HTML}{e14775}   % keywords (pink)
\definecolor{lpModifier}{HTML}{4a7fb5}  % modifiers (muted blue)
\definecolor{lpType}{HTML}{7058be}      % types (purple)
\definecolor{lpComment}{HTML}{3b5998}   % comments (dark blue)
\definecolor{lpPunct}{HTML}{918c8e}     % punctuation (grey)
\definecolor{lpFg}{HTML}{3a3338}        % foreground
\definecolor{lpBoxBg}{HTML}{fbfaf9}     % box background
\definecolor{lpBoxFr}{HTML}{c5c1bf}     % box frame

% ── Math-mode shorthand ──────────────────────────────────────────
\newcommand{\lmark}[1]{\textcolor{lpPunct}{#1}}

% ── LambdaPi listings language ────────────────────────────────────
\lstdefinelanguage{lambdapi}{
  comment=[l]{//},
  alsoletter={_},
  keywords={symbol, rule, with, assert, require, open, type,
            begin, end, assume, refine, have, rewrite,
            admit, admitted, apply, simplify, reflexivity, symmetry},
  keywords=[2]{constant, sequential, opaque, injective},
  keywords=[3]{TYPE, Prop, Set},
  sensitive=true,
}

\lstdefinestyle{lp}{
  language=lambdapi,
  basicstyle=\footnotesize\ttfamily\color{lpFg},
  keywordstyle=\color{lpKeyword}\bfseries,
  keywordstyle=[2]\color{lpModifier},
  keywordstyle=[3]\color{lpType},
  commentstyle=\color{lpComment}\itshape,
  showstringspaces=false,
  columns=flexible,
  keepspaces=true,
  mathescape=true,
  frame=none,
  numbers=none,
  literate=
    {(}{{\lmark{(}}}1
    {)}{{\lmark{)}}}1
    {[}{{\lmark{[}}}1
    {]}{{\lmark{]}}}1
    {;}{{\lmark{;}}}1
    {:}{{\lmark{:}}}1,
  aboveskip=0pt,
  belowskip=0pt,
}

% ── PP replay listing style ──────────────────────────────────────
\lstdefinestyle{ppreplay}{
  basicstyle=\footnotesize\ttfamily\color{lpFg},
  showstringspaces=false,
  columns=flexible,
  keepspaces=true,
  mathescape=true,
  frame=none,
  numbers=none,
  literate=
    {[}{{\lmark{[}}}1
    {]}{{\lmark{]}}}1,
  aboveskip=0pt,
  belowskip=0pt,
}

\pagestyle{empty}
\begin{document}

\paragraph{Verification of Predicate Prover proofs in LambdaPi.}
%
The Predicate Prover (PP) is an automated theorem prover
used by Atelier~B to discharge proof obligations arising from
B-method developments.
%
Because PP's source code is not publicly available,
its results must be trusted without independent verification.
%
We address this gap with \texttt{pp2lp}, an OCaml tool that
reconstructs PP proofs in LambdaPi,
a proof assistant based on the $\lambda\Pi$-calculus modulo rewriting.

We first encode PP's inference rules
(covering conjunction, disjunction, implication,
negation, quantification, equality, normalisation,
arithmetic, booleans)
as LambdaPi symbols whose types capture
each rule's antecedents and consequent, and whose bodies
are proofs of soundness (\autoref{fig:pp-rule}).
%
PP can be instrumented to record a compact
\emph{trace} listing the inference rules applied
during a proof.
%
A companion Atelier~B utility then expands this trace
into a \emph{replay} that records, for each step,
the rule applied and the intermediate goal
(\autoref{fig:pp-trace}, left).
%
Given a replay, \texttt{pp2lp} parses it,
reconstructs a proof tree, and emits a LambdaPi script
that composes the encoded rules following the tree structure,
yielding a complete, independently checkable proof
(\autoref{fig:pp-trace}, right).
%
We currently encode 126 of PP's 140~inference rules
and successfully verify over 100~proof goals from a
custom benchmark suite.

% ----- Figure 1: PP rule -> LambdaPi encoding -----
\begin{figure}[ht]
	\centering
	\begin{tikzpicture}[
			lbl/.style={font=\footnotesize\bfseries, text=lpFg},
			arr/.style={->, thick, >=stealth, color=lpPunct},
		]
		% PP specification (left) — tabular matching spec_pp style
		\node[anchor=east, inner sep=0pt] (spec) at (-0.5,0) {%
			\begin{tcolorbox}[
					colback=lpBoxBg, colframe=lpBoxFr,
					arc=2pt, boxrule=0.4pt,
					left=5pt, right=5pt, top=3pt, bottom=3pt,
					width=6.2cm,
				]
				\centering\footnotesize
				\begin{tabular}{@{}l@{\quad}c@{\quad}c@{}}
					\toprule
					                & \textbf{Ant\'{e}c\'{e}dents}              & \textbf{Cons\'{e}quent} \\
					\midrule
					$\mathsf{AND2}$ & $\mathsf{H} \vdash P \Rightarrow \lnot Q$
					                & $\mathsf{H} \vdash \lnot(P \land Q)$                                \\
					\bottomrule
				\end{tabular}
			\end{tcolorbox}};
		\node[lbl, anchor=south west] at (spec.north west) {\textbf{PP specification}};

		% LambdaPi encoding (right)
		\node[anchor=west, inner sep=0pt] (lp) at (0.5,0) {%
			\begin{tcolorbox}[
					colback=lpBoxBg, colframe=lpBoxFr,
					arc=2pt, boxrule=0.4pt,
					left=5pt, right=5pt, top=3pt, bottom=3pt,
					width=5.6cm,
				]
				\begin{lstlisting}[style=lp]
opaque symbol AND2 [P Q]
  : $\pi$ (P $\Rightarrow$ $\lnot$ Q)
  $\to$ $\pi$ ($\lnot$ (P $\land$ Q)) $\coloneqq$
begin
  assume P Q h1 h2;
  refine h1 (${\land}_{e_1}$ h2) (${\land}_{e_2}$ h2)
end;
\end{lstlisting}
			\end{tcolorbox}};
		\node[lbl, anchor=south west] at (lp.north west) {\textbf{LambdaPi encoding \& proof}};

		% Arrow
		\draw[arr] (spec.east) -- (lp.west);
	\end{tikzpicture}
	\caption{PP rule $\mathsf{AND2}$ and its LambdaPi encoding.
		The antecedent becomes a function parameter and the consequent
		the return type; the proof body witnesses soundness
		using conjunction eliminators ${\land}_{e_1}$, ${\land}_{e_2}$
		from LambdaPi's standard library.}
	\label{fig:pp-rule}
\end{figure}

% ----- Figure 2: PP replay -> LambdaPi proof -----
\begin{figure}[ht]
	\centering
	\begin{tikzpicture}[
			lbl/.style={font=\footnotesize\bfseries, text=lpFg},
			arr/.style={->, thick, >=stealth, color=lpPunct},
		]
		% PP replay (left)
		\node[anchor=east, inner sep=0pt] (trace) at (-0.5,0) {%
			\begin{tcolorbox}[
					colback=lpBoxBg, colframe=lpBoxFr,
					arc=2pt, boxrule=0.4pt,
					left=5pt, right=5pt, top=3pt, bottom=3pt,
					width=5.0cm,
				]
				\begin{lstlisting}[style=ppreplay]
[$\textbf{AND2}$] <not(p and not(p))>
[$\textbf{IMP4}$] <p => not(not(p))>
[$\textbf{NOT2}$] <not(not(p))>
[$\textbf{AXM3}$] <p>
\end{lstlisting}
			\end{tcolorbox}};
		\node[lbl, anchor=south west] at (trace.north west) {\textbf{PP proof replay}};

		% LambdaPi proof (right)
		\node[anchor=west, inner sep=0pt] (lp) at (0.5,0) {%
			\begin{tcolorbox}[
					colback=lpBoxBg, colframe=lpBoxFr,
					arc=2pt, boxrule=0.4pt,
					left=5pt, right=5pt, top=3pt, bottom=3pt,
					width=5.6cm,
				]
				\begin{lstlisting}[style=lp]
opaque symbol t02 (p : Prop)
  : $\pi$ ($\lnot$ (p $\land$ $\lnot$ p)) $\coloneqq$
begin
  assume p;
  refine AND2 _;
  refine IMP4 _;
  assume hp;
  refine NOT2 hp
end;
\end{lstlisting}
			\end{tcolorbox}};
		\node[lbl, anchor=south west] at (lp.north west) {\textbf{LambdaPi reconstruction}};

		% Arrow
		\draw[arr] (trace.east) -- (lp.west);
	\end{tikzpicture}
	\caption{Automated reconstruction of a PP proof of $\lnot(p \land \lnot p)$.
		Each line in the replay maps to a \texttt{refine} tactic
		that applies the corresponding rule;
		LambdaPi type-checks the composed proof against the goal type.}
	\label{fig:pp-trace}
\end{figure}

\end{document}
